% !TEX root = ../main.tex
%% Chapitre 6 — Fonctions trigonométriques
%% Sources : frontend/cours/resume/terminal/resume_fonctions_trigonometriques_terminal.txt
%%           frontend/exercices/terminal_spé_Maths/exercice_fonctions_trigonometriques.txt
%% L'exercice « Intégration et applications » de la source est déplacé au chapitre 9 (Calcul intégral).
%% Exercices 2, 13, 14 et 15 : exercices inédits rédigés pour le livre (numérotation de l'édition 2027-2028).
%% Édition 2027-2028 : section 2 de L'essentiel réécrite ; exercices 3 et 4 nouveaux (programme de la rentrée 2027).

\begin{ouverture}[trigo]{6}{Fonctions trigonométriques}{Partie 1 · Analyse}
  {Cercle trigonométrique\enspace·\enspace Sinus et cosinus\enspace·\enspace Dérivées\enspace·\enspace
   Équations\enspace·\enspace Inéquations}
  \objectif{Lire sur le cercle trigonométrique les valeurs remarquables et utiliser les angles associés.}
  \objectif{Résoudre une équation ou une inéquation trigonométrique sur un intervalle donné.}
  \objectif{Dériver une fonction construite avec $\sin$ et $\cos$, et étudier ses variations.}
  \objectif{Relier le cercle trigonométrique aux courbes de $\sin$ et $\cos$ ; traduire graphiquement parité et périodicité.}
  \objectif{Connaître et utiliser la limite de $\dfrac{\sin x}{x}$ en $0$.}
  \prerequis{Cercle trigonométrique, radian, cosinus et sinus d'un réel, valeurs remarquables (Première).}
  \prerequis{Fonctions paires et impaires (Première) ; dérivation et fonctions composées (chapitre 3).}
  \prerequis{Limites, continuité et théorème des valeurs intermédiaires (chapitres 2 et 4).}
\end{ouverture}

%% =====================================================================
%%  L'ESSENTIEL
%% =====================================================================
\essentiel

\essSection{1}{Le cercle trigonométrique}

\begin{definition}[Cosinus et sinus d'un réel]
Le plan est muni d'un repère orthonormé $(O\pv I,J)$. Le \textbf{cercle trigonométrique} est le cercle de centre $O$ et de
rayon $1$, parcouru dans le sens direct. À tout réel $x$ correspond un point $M(x)$ du cercle : $x$ est une mesure en
radians de l'angle orienté $\big(\overrightarrow{OI},\overrightarrow{OM}\big)$. Le \textbf{cosinus} et le \textbf{sinus} de $x$
sont les coordonnées de ce point : $M(x)\,(\cos x\pv\sin x)$.
\end{definition}

\noindent
\begin{minipage}[c]{0.4\linewidth}
\centering
\begin{tikzpicture}
  \coordinate (O) at (0,0);
  \coordinate (M) at (50:2);
  \draw[otMuted,line width=0.5pt,-{Stealth[length=2mm]}] (-2.5,0) -- (2.6,0);
  \draw[otMuted,line width=0.5pt,-{Stealth[length=2mm]}] (0,-2.5) -- (0,2.6);
  \draw[otNavy,line width=1.1pt] (O) circle (2);
  \draw[otGreen,line width=0.6pt,dash pattern=on 2pt off 1.5pt] (M) -- (M |- O);
  \draw[otGreen,line width=0.6pt,dash pattern=on 2pt off 1.5pt] (M) -- (O |- M);
  \draw[otNavy,line width=0.7pt] (O) -- (M);
  \draw[otGreenDark,line width=0.7pt,-{Stealth[length=1.5mm]}] (0.55,0) arc (0:50:0.55);
  \node[font=\scriptsize,text=otGreenDark] at (24:0.82) {$x$};
  \fill[otNavy] (M) circle (1.7pt);
  \node[anchor=south west,font=\small,inner sep=1.5pt] at (M) {$M(x)$};
  \fill[otGreenDark] (M |- O) circle (1.3pt);
  \fill[otGreenDark] (O |- M) circle (1.3pt);
  \node[anchor=north,font=\scriptsize,text=otGreenDark] at ([yshift=-1pt]M |- O) {$\cos x$};
  \node[anchor=east,font=\scriptsize,text=otGreenDark] at ([xshift=-1pt]O |- M) {$\sin x$};
  \node[anchor=north east,font=\scriptsize,inner sep=1.5pt] at (O) {$O$};
  \node[anchor=north west,font=\scriptsize,inner sep=1.5pt] at (2,0) {$I$};
  \node[anchor=south east,font=\scriptsize,inner sep=1.5pt] at (0,2) {$J$};
\end{tikzpicture}
\end{minipage}\hfill
\begin{minipage}[c]{0.57\linewidth}
\centering
{\renewcommand{\arraystretch}{2}\arrayrulecolor{otRule}%
\begin{tabular}{@{}>{\leavevmode\color{otNavy}}l*{6}{>{\centering\arraybackslash}p{8.4mm}}@{}}
\toprule
$x$ & $0$ & $\dfrac{\pi}{6}$ & $\dfrac{\pi}{4}$ & $\dfrac{\pi}{3}$ & $\dfrac{\pi}{2}$ & $\pi$ \\
\midrule
$\cos x$ & $1$ & $\dfrac{\sqrt3}{2}$ & $\dfrac{\sqrt2}{2}$ & $\dfrac12$ & $0$ & $-1$ \\
$\sin x$ & $0$ & $\dfrac12$ & $\dfrac{\sqrt2}{2}$ & $\dfrac{\sqrt3}{2}$ & $1$ & $0$ \\
\bottomrule
\end{tabular}}
\par\vspace{2mm}
\legende{Valeurs remarquables, à connaître par cœur.}
\end{minipage}

\begin{propriete}[Pour tout réel $x$ et tout entier $k$]
\centering
$-1\leq\cos x\leq1$ \qquad $-1\leq\sin x\leq1$ \qquad $\cos^{2}x+\sin^{2}x=1$\\[1.2mm]
$\cos(x+2k\pi)=\cos x$ \qquad $\sin(x+2k\pi)=\sin x$
\end{propriete}

\par\vspace{2mm}
\noindent
\begin{minipage}[c]{0.4\linewidth}
\centering
\begin{tikzpicture}
  \coordinate (O) at (0,0);
  \draw[otMuted,line width=0.5pt,-{Stealth[length=2mm]}] (-2.4,0) -- (2.5,0);
  \draw[otMuted,line width=0.5pt,-{Stealth[length=2mm]}] (0,-2.4) -- (0,2.5);
  \draw[otMuted,line width=0.5pt,dash pattern=on 3pt off 2pt] (-1.5,-1.5) -- (1.9,1.9);
  \draw[otNavy,line width=1.1pt] (O) circle (1.8);
  \draw[otGreen,line width=0.6pt,dash pattern=on 2pt off 1.5pt] (25:1.8) -- (-25:1.8);
  \draw[otGreen,line width=0.6pt,dash pattern=on 2pt off 1.5pt] (25:1.8) -- (155:1.8);
  \draw[otGreen,line width=0.6pt,dash pattern=on 2pt off 1.5pt] (25:1.8) -- (205:1.8);
  \draw[otGreen,line width=0.6pt,dash pattern=on 2pt off 1.5pt] (25:1.8) -- (65:1.8);
  \fill[otNavy] (25:1.8) circle (1.7pt);
  \foreach \a in {-25,155,205,65} \fill[otGreenDark] (\a:1.8) circle (1.4pt);
  \node[anchor=west,font=\scriptsize] at (25:1.86) {$x$};
  \node[anchor=north west,font=\scriptsize] at (-25:1.84) {$-x$};
  \node[anchor=south east,font=\scriptsize] at (155:1.84) {$\pi-x$};
  \node[anchor=north east,font=\scriptsize] at (205:1.84) {$\pi+x$};
  \node[anchor=south west,font=\scriptsize] at (65:1.84) {$\frac{\pi}{2}-x$};
\end{tikzpicture}
\end{minipage}\hfill
\begin{minipage}[c]{0.57\linewidth}
\raggedright
Les points associés à $x$ se déduisent de $M(x)$ par symétrie : par rapport à l'axe des abscisses pour $-x$, à l'axe
des ordonnées pour $\pi-x$, à l'origine pour $\pi+x$ et à la droite $y=x$ pour $\frac{\pi}{2}-x$.
\par\vspace{3mm}\centering
{\renewcommand{\arraystretch}{1.75}\arrayrulecolor{otRule}%
\begin{tabular}{@{}>{\leavevmode\color{otNavy}}l>{\centering\arraybackslash}p{24mm}>{\centering\arraybackslash}p{24mm}@{}}
\toprule
 & $\cos$ & $\sin$ \\
\midrule
$-x$ & $\cos x$ & $-\sin x$ \\
$\pi-x$ & $-\cos x$ & $\sin x$ \\
$\pi+x$ & $-\cos x$ & $-\sin x$ \\
$\dfrac{\pi}{2}-x$ & $\sin x$ & $\cos x$ \\
\bottomrule
\end{tabular}}
\end{minipage}

\essSection{2}{Les fonctions sinus et cosinus}

\begin{definition}[Fonctions cosinus et sinus\enspace\otNouveau]
La fonction \textbf{cosinus} associe à tout réel $x$ le réel $\cos x$, abscisse du point $M(x)$ du cercle trigonométrique ; la
fonction \textbf{sinus} lui associe $\sin x$, ordonnée de $M(x)$. Ces deux fonctions sont définies sur $\R$ et prennent leurs
valeurs dans $\intFF{-1}{1}$.
\end{definition}

\noindent
\begin{minipage}[c]{0.66\linewidth}
\centering
\begin{tikzpicture}[x=1cm,y=1cm]
  \draw[otMuted,line width=0.5pt,-{Stealth[length=1.6mm]}] (-1.35,0) -- (1.4,0);
  \draw[otMuted,line width=0.5pt,-{Stealth[length=1.6mm]}] (0,-1.3) -- (0,1.45);
  \draw[otNavy,line width=1pt] (0,0) circle (1);
  \draw[otGreen,line width=2.2pt] (1,0) arc (0:120:1);
  \coordinate (M) at (120:1);
  \draw[otNavy,line width=0.6pt] (0,0) -- (M);
  \node[anchor=north east,font=\scriptsize,inner sep=1.5pt] at (0,0) {$O$};
  \node[anchor=north west,font=\scriptsize,inner sep=1.5pt] at (1,0) {$1$};
  \begin{scope}[shift={(2.1,0)}]
    \draw[otMuted,line width=0.5pt,-{Stealth[length=1.6mm]}] (-0.2,0) -- (6.75,0);
    \draw[otMuted,line width=0.5pt,-{Stealth[length=1.6mm]}] (0,-1.3) -- (0,1.45);
    \foreach \xv/\lab in {1.5708/{\frac{\pi}{2}},3.1416/{\pi},4.7124/{\frac{3\pi}{2}},6.2832/{2\pi}}{%
      \draw[otMuted,line width=0.4pt] (\xv,-0.05) -- (\xv,0.05);
      \node[anchor=north,font=\scriptsize,text=otMuted] at (\xv,-0.06) {$\lab$};}
    \draw[otMuted,line width=0.4pt] (-0.05,1) -- (0.05,1) node[anchor=east,font=\scriptsize,text=otMuted,xshift=-1.5pt] {$1$};
    \draw[otMuted,line width=0.4pt] (-0.05,-1) -- (0.05,-1) node[anchor=east,font=\scriptsize,text=otMuted,xshift=-1.5pt] {$-1$};
    \draw[otNavy,line width=1.1pt,domain=0:6.2832,samples=150] plot (\x,{sin(\x r)});
    \draw[otGreen,line width=2.2pt] (0,0) -- (2.0944,0);
    \coordinate (P) at (2.0944,0.8660);
    \draw[otGreenDark,line width=0.5pt,dash pattern=on 2pt off 1.5pt] (P) -- (2.0944,0);
    \node[anchor=north,font=\scriptsize,text=otGreenDark] at (2.0944,-0.05) {$x$};
    \node[anchor=west,font=\scriptsize,text=otNavy] at (5.05,-1.05) {$y=\sin x$};
  \end{scope}
  \draw[otGreenDark,line width=0.6pt,dash pattern=on 2pt off 1.5pt] (M) -- (P);
  \fill[otNavy] (M) circle (1.7pt);
  \node[anchor=south east,font=\scriptsize,inner sep=1.5pt] at (M) {$M(x)$};
  \fill[otGreenDark] (P) circle (1.7pt);
  \node[anchor=south west,font=\scriptsize,text=otGreenDark,inner sep=1.5pt] at (P) {$(x\pv\sin x)$};
\end{tikzpicture}
\end{minipage}\hfill
\begin{minipage}[c]{0.31\linewidth}
\raggedright\fontsize{10}{13.5}\selectfont
On enroule la droite des réels sur le cercle : l'arc vert mesure $x$. L'ordonnée de $M(x)$, reportée au-dessus de
l'abscisse $x$, donne un point de la courbe de sinus. On construit de même la courbe de cosinus avec l'abscisse de $M(x)$.
\end{minipage}

\begin{propriete}[Parité et périodicité]
Pour tout réel $x$ : $\cos(x+2\pi)=\cos x$, $\sin(x+2\pi)=\sin x$, $\cos(-x)=\cos x$ et $\sin(-x)=-\sin x$. Les fonctions $\cos$ et
$\sin$ sont donc \textbf{périodiques de période $2\pi$} : leurs courbes sont invariantes par la translation de vecteur
$2\pi\,\vec\imath$. La fonction $\cos$ est \textbf{paire} (courbe symétrique par rapport à l'axe des ordonnées) et la fonction
$\sin$ est \textbf{impaire} (courbe symétrique par rapport à l'origine).
\end{propriete}

\begin{demonstration}[lecture sur le cercle]
Un tour complet du cercle mesure $2\pi$ : les réels $x$ et $x+2\pi$ ont le même point image, donc le même cosinus et le même
sinus. Le point $M(-x)$ est le symétrique de $M(x)$ par rapport à l'axe des abscisses : même abscisse, ordonnées opposées.
\end{demonstration}

\begin{center}
\begin{tikzpicture}
\begin{axis}[otGraphe, width=\linewidth, xmin=-7, xmax=7.3, ymin=-1.35, ymax=1.6,
  xtick={-6.2832,-4.7124,-3.1416,-1.5708,1.5708,3.1416,4.7124,6.2832},
  xticklabels={$-2\pi$,$-\frac{3\pi}{2}$,$-\pi$,$-\frac{\pi}{2}$,$\frac{\pi}{2}$,$\pi$,$\frac{3\pi}{2}$,$2\pi$},
  ytick={-1,1}]
  \addplot[otGreen,line width=1.2pt,domain=-7:7,samples=240] {cos(deg(x))};
  \addplot[otNavy,line width=1.2pt,domain=-7:7,samples=240] {sin(deg(x))};
  \node[anchor=south,otEtiquette,font=\scriptsize,text=otNavy] at (axis cs:2.2,1.05) {$y=\sin x$};
  \node[anchor=south,otEtiquette,font=\scriptsize,text=otGreenDark] at (axis cs:-0.75,1.05) {$y=\cos x$};
\end{axis}
\end{tikzpicture}\\[0.5mm]
\legende{La courbe de $\cos$ se déduit de celle de $\sin$ par la translation de vecteur $-\frac{\pi}{2}\,\vec\imath$, car $\cos x=\sin\left(x+\frac{\pi}{2}\right)$.}
\end{center}

\begin{theoreme}[Dérivées]
Les fonctions $\sin$ et $\cos$ sont dérivables sur $\R$ (admis) et $\sin'=\cos$, $\cos'=-\sin$. Plus généralement, si $u$
est dérivable sur un intervalle $I$ :
\[ \big(\sin u\big)'=u'\cos u \qquad\qquad \big(\cos u\big)'=-u'\sin u \qquad\qquad
   \text{par exemple}\enspace \big(\cos(ax+b)\big)'=-a\sin(ax+b). \]
\end{theoreme}

\noindent
\begin{minipage}[t]{0.48\linewidth}
\centering
\begin{tableau}
  \otTab[1.5]{$x$/0.8,$-\sin x$/0.8,$\cos$/1.6}{$0$,$\frac{\pi}{2}$,$\pi$}
  \tkzTabLine{z,-,t,-,z}
  \tkzTabVar{+/$1$,R/,-/$-1$}
  \tkzTabIma{1}{3}{2}{$0$}
\end{tableau}
\end{minipage}\hfill
\begin{minipage}[t]{0.48\linewidth}
\centering
\begin{tableau}
  \otTab[1.5]{$x$/0.8,$\cos x$/0.8,$\sin$/1.6}{$0$,$\frac{\pi}{2}$,$\pi$}
  \tkzTabLine{,+,z,-,}
  \tkzTabVar{-/$0$,+/$1$,-/$0$}
\end{tableau}
\end{minipage}
\par\vspace{1mm}
\begin{center}\legende{Variations sur $\intFF{0}{\pi}$ ; on complète par parité puis par périodicité.}\end{center}

\begin{theoreme}[Une limite de référence]
\centering $\displaystyle\lim_{x\to0}\frac{\sin x}{x}=1$
\end{theoreme}

\begin{demonstration}
Pour $x\neq0$, $\dfrac{\sin x}{x}=\dfrac{\sin x-\sin0}{x-0}$ est le taux d'accroissement de $\sin$ entre $0$ et $x$. Comme $\sin$
est dérivable en $0$, il tend vers $\sin'(0)=\cos0=1$ lorsque $x$ tend vers $0$.
\end{demonstration}

\begin{savoirfaire}[Étudier une fonction trigonométrique]
\begin{enumerate}[leftmargin=1.4em,itemsep=0.4mm,topsep=0.5mm,label={\color{otNavy}\arabic*.}]
  \item Réduire l'intervalle d'étude : une fonction $T$-périodique s'étudie sur un intervalle de longueur $T$ ; si elle est
    de plus paire ou impaire, sur la moitié de cet intervalle.
  \item Dériver, puis \textbf{factoriser} $f'(x)$, souvent à l'aide de $\cos x$ ou de $\sin x$.
  \item Étudier le signe de chaque facteur en lisant le cercle trigonométrique.
  \item Compléter le tableau et la courbe par symétrie, puis par translation.
\end{enumerate}
\end{savoirfaire}

\essSection{3}{Équations et inéquations}

\begin{savoirfaire}[Résoudre une équation trigonométrique]
Pour tous réels $x$ et $a$, avec $k$ décrivant $\mathbb{Z}$ :
\[ \cos x=\cos a\iff x=a+2k\pi\ \text{ ou }\ x=-a+2k\pi \qquad
   \sin x=\sin a\iff x=a+2k\pi\ \text{ ou }\ x=\pi-a+2k\pi \]
On écrit d'abord la constante comme le cosinus ou le sinus d'un angle remarquable, puis on ne garde que les solutions
de l'intervalle demandé.
\end{savoirfaire}

\noindent
\begin{minipage}[t]{0.32\linewidth}
\centering
\begin{tikzpicture}[scale=0.82]
  \draw[otMuted,line width=0.5pt,-{Stealth[length=1.6mm]}] (-1.55,0) -- (1.65,0);
  \draw[otMuted,line width=0.5pt,-{Stealth[length=1.6mm]}] (0,-1.55) -- (0,1.65);
  \draw[otNavy,line width=1pt] (0,0) circle (1.25);
  \draw[otGreen,line width=0.7pt,dash pattern=on 2pt off 1.5pt] (50:1.25) -- (-50:1.25);
  \fill[otNavy] (50:1.25) circle (1.5pt);
  \fill[otNavy] (-50:1.25) circle (1.5pt);
  \fill[otGreenDark] (0.8035,0) circle (1.2pt);
  \node[anchor=south west,font=\scriptsize,inner sep=1.5pt] at (50:1.25) {$a$};
  \node[anchor=north west,font=\scriptsize,inner sep=1.5pt] at (-50:1.25) {$-a$};
  \node[anchor=north east,font=\scriptsize,text=otGreenDark,inner sep=1.5pt] at (0.8035,0) {$\cos a$};
\end{tikzpicture}\\[1mm]
\legende{$\cos x=\cos a$}
\end{minipage}\hfill
\begin{minipage}[t]{0.32\linewidth}
\centering
\begin{tikzpicture}[scale=0.82]
  \draw[otMuted,line width=0.5pt,-{Stealth[length=1.6mm]}] (-1.55,0) -- (1.65,0);
  \draw[otMuted,line width=0.5pt,-{Stealth[length=1.6mm]}] (0,-1.55) -- (0,1.65);
  \draw[otNavy,line width=1pt] (0,0) circle (1.25);
  \draw[otGreen,line width=0.7pt,dash pattern=on 2pt off 1.5pt] (35:1.25) -- (145:1.25);
  \fill[otNavy] (35:1.25) circle (1.5pt);
  \fill[otNavy] (145:1.25) circle (1.5pt);
  \fill[otGreenDark] (0,0.717) circle (1.2pt);
  \node[anchor=south west,font=\scriptsize,inner sep=1.5pt] at (35:1.25) {$a$};
  \node[anchor=south east,font=\scriptsize,inner sep=1.5pt] at (145:1.25) {$\pi-a$};
  \node[anchor=south west,font=\scriptsize,text=otGreenDark,inner sep=1.5pt] at (0,0.75) {$\sin a$};
\end{tikzpicture}\\[1mm]
\legende{$\sin x=\sin a$}
\end{minipage}\hfill
\begin{minipage}[t]{0.32\linewidth}
\centering
\begin{tikzpicture}[scale=0.82]
  \draw[otMuted,line width=0.5pt,-{Stealth[length=1.6mm]}] (-1.55,0) -- (1.65,0);
  \draw[otMuted,line width=0.5pt,-{Stealth[length=1.6mm]}] (0,-1.55) -- (0,1.65);
  \draw[otNavy,line width=1pt] (0,0) circle (1.25);
  \draw[otGreen,line width=2.4pt] (-60:1.25) arc (-60:60:1.25);
  \draw[otGreen,line width=0.7pt,dash pattern=on 2pt off 1.5pt] (0.625,-1.4) -- (0.625,1.4);
  \fill[otGreenDark] (60:1.25) circle (1.5pt);
  \fill[otGreenDark] (-60:1.25) circle (1.5pt);
  \node[anchor=south west,font=\scriptsize,inner sep=1.5pt] at (60:1.25) {$\frac{\pi}{3}$};
  \node[anchor=north west,font=\scriptsize,inner sep=1.5pt] at (-60:1.25) {$-\frac{\pi}{3}$};
  \node[anchor=north east,font=\scriptsize,text=otGreenDark,inner sep=1.5pt] at (0.625,0) {$\frac12$};
\end{tikzpicture}\\[1mm]
\legende{$\cos x\geq\frac12$ sur $\intOF{-\pi}{\pi}$ : $x\in\intFF{-\frac{\pi}{3}}{\frac{\pi}{3}}$}
\end{minipage}

\begin{savoirfaire}[Résoudre une inéquation trigonométrique]
On résout d'abord l'équation associée sur l'intervalle demandé, puis on lit sur le cercle l'arc formé des points dont
l'abscisse (pour $\cos$) ou l'ordonnée (pour $\sin$) convient. Pour une inéquation en $\cos(2x)$ avec $x\in\intFF{0}{\pi}$,
on pose $X=2x$, qui décrit $\intFF{0}{2\pi}$.
\end{savoirfaire}

\essSection{4}{Pour aller plus loin}

\begin{propriete}[Formules d'addition et de duplication]
Pour tous réels $a$ et $b$ :\enspace
$\cos(a+b)=\cos a\cos b-\sin a\sin b$\qquad $\sin(a+b)=\sin a\cos b+\cos a\sin b$\\[1mm]
$\cos(2a)=\cos^{2}a-\sin^{2}a=2\cos^{2}a-1=1-2\sin^{2}a$\qquad $\sin(2a)=2\sin a\cos a$\\[1mm]
{\small\color{otMuted}Ces formules ne figurent pas au programme de spécialité : lorsqu'elles sont utiles, l'énoncé les rappelle.}
\end{propriete}

\noindent
\begin{minipage}[c]{0.55\linewidth}
\raggedright
\begin{definition}[La fonction tangente]
\raggedright
Pour $x\neq\frac{\pi}{2}+k\pi$ ($k\in\mathbb{Z}$), on pose $\tan x=\dfrac{\sin x}{\cos x}$. La fonction $\tan$ est impaire,
périodique de période $\pi$, et
\[ \tan'(x)=\frac{1}{\cos^{2}x}=1+\tan^{2}x. \]
Elle est strictement croissante sur $\intOO{-\frac{\pi}{2}}{\frac{\pi}{2}}$, avec $\tan\frac{\pi}{4}=1$.
\end{definition}
{\small\color{otMuted}La fonction tangente n'est pas au programme de spécialité ; elle apparaît dans quelques exercices.}
\end{minipage}\hfill
\begin{minipage}[c]{0.41\linewidth}
\centering
\begin{tikzpicture}
\begin{axis}[otGraphe, width=58mm, xmin=-4.2, xmax=4.4, ymin=-3.2, ymax=3.4,
  xtick={-3.1416,-1.5708,1.5708,3.1416}, xticklabels={$-\pi$,$-\frac{\pi}{2}$,$\frac{\pi}{2}$,$\pi$},
  ytick={-2,-1,1,2}]
  \draw[otGreen,dashed,line width=0.6pt] (axis cs:1.5708,-3.2) -- (axis cs:1.5708,3.3);
  \draw[otGreen,dashed,line width=0.6pt] (axis cs:-1.5708,-3.2) -- (axis cs:-1.5708,3.3);
  \addplot[otNavy,line width=1.2pt,domain=-1.26:1.26,samples=100] {tan(deg(x))};
  \addplot[otNavy,line width=1.2pt,domain=1.885:4.2,samples=100] {tan(deg(x))};
  \addplot[otNavy,line width=1.2pt,domain=-4.2:-1.885,samples=100] {tan(deg(x))};
\end{axis}
\end{tikzpicture}\\[0.5mm]
\legende{Les droites $x=\frac{\pi}{2}+k\pi$ sont asymptotes.}
\end{minipage}

\begin{attention}[Pièges fréquents]
\begin{itemize}[label={\color{otOrange}\textbullet},leftmargin=1.4em,itemsep=0.6mm,topsep=0.5mm]
  \item $\cos^{2}x+\sin^{2}x=1$, mais en général $\cos x+\sin x\neq1$.
  \item $(\cos x)'=-\sin x$ : attention au signe. Et $\big(\sin(ax+b)\big)'=a\cos(ax+b)$ : ne pas oublier le facteur $a$.
  \item Une équation trigonométrique a en général une infinité de solutions : écrire « $+2k\pi$, $k\in\mathbb{Z}$ », puis ne
    garder que les solutions de l'intervalle demandé.
  \item Si $x$ décrit $\intFF{0}{\pi}$, alors $2x$ décrit $\intFF{0}{2\pi}$ : on change d'intervalle en même temps que de variable.
  \item La calculatrice doit être réglée en \textbf{radians}.
\end{itemize}
\end{attention}

%% =====================================================================
%%  EXERCICES
%% =====================================================================
\exercices
\rubrique{Découverte}

\begin{exercice}[1]{Valeurs remarquables}
Calculer la valeur exacte de chacune des expressions suivantes.
\begin{tasks}(3)
  \task $\sin\dfrac{\pi}{6}+\cos\dfrac{\pi}{3}$
  \task $\tan\dfrac{\pi}{4}-\sin\dfrac{\pi}{2}$
  \task $\sin^{2}\dfrac{\pi}{4}+\cos^{2}\dfrac{\pi}{4}$
  \task $\dfrac{\sin\frac{\pi}{6}}{\cos\frac{\pi}{6}}$
  \task $\sin\dfrac{\pi}{3}\times\cos\dfrac{\pi}{6}$
\end{tasks}
\end{exercice}

\begin{exercice}[1]{Angles associés}
\begin{questions}
  \item Calculer les valeurs exactes suivantes, en se ramenant à un angle de $\intFF{0}{\frac{\pi}{2}}$.
\begin{tasks}(3)
  \task $\cos\dfrac{5\pi}{6}$
  \task $\sin\dfrac{7\pi}{6}$
  \task $\cos\left(-\dfrac{2\pi}{3}\right)$
  \task $\sin\dfrac{3\pi}{4}$
  \task $\cos\dfrac{11\pi}{6}$
  \task $\sin\dfrac{17\pi}{4}$
\end{tasks}
  \item Simplifier, pour tout réel $x$, les expressions suivantes.
    \begin{gather*}
      A(x)=\cos(\pi-x)+\cos(-x) \qquad B(x)=\sin(\pi-x)-\sin(-x) \\
      C(x)=\sin\left(\tfrac{\pi}{2}-x\right)+\cos(\pi+x)+\cos(x+2\pi)
    \end{gather*}
  \item Sachant que $\cos x=\dfrac35$ et que $x\in\intFF{-\frac{\pi}{2}}{0}$, calculer $\sin x$.
\end{questions}
\end{exercice}

\begin{exercice}[1]{Du cercle à la courbe\enspace\otNouveau}
\begin{codecote}[0.57]
\begin{questions}
  \item Placer sur un cercle trigonométrique les points $M\left(\frac{\pi}{6}\right)$, $M\left(\frac{5\pi}{6}\right)$ et
    $M\left(-\frac{\pi}{6}\right)$. En déduire $\sin\frac{\pi}{6}$, $\sin\frac{5\pi}{6}$ et $\sin\left(-\frac{\pi}{6}\right)$.
  \item Expliquer à l'aide du cercle pourquoi $\sin(\pi-x)=\sin x$ pour tout réel $x$. Quelle symétrie de la courbe de la
    fonction sinus cette égalité traduit-elle ?
  \item Lorsque $x$ croît de $0$ à $\frac{\pi}{2}$, puis de $\frac{\pi}{2}$ à $\pi$, comment évolue l'ordonnée de $M(x)$ ?
    En déduire les variations de la fonction sinus sur $\intFF{0}{\pi}$.
\end{questions}
\cote
\centering
\begin{tikzpicture}
\begin{axis}[otGraphe, unit vector ratio=1 1.6, width=64mm, xmin=-3.6, xmax=3.75, ymin=-1.35, ymax=1.45,
  xtick={-3.1416,-1.0472,1.0472,3.1416}, xticklabels={$-\pi$,$-\frac{\pi}{3}$,$\frac{\pi}{3}$,$\pi$},
  ytick={-1,0.5,1}, yticklabels={$-1$,$\frac12$,$1$}, ylabel={}]
  \draw[otOrange,line width=0.7pt,dash pattern=on 3pt off 2pt] (axis cs:-3.45,0.5) -- (axis cs:3.45,0.5);
  \addplot[otGreen,line width=1.2pt,domain=-3.1416:3.1416,samples=160] {cos(deg(x))};
  \node[anchor=south,font=\scriptsize,text=otGreenDark] at (axis cs:2.35,-1.3) {$y=\cos x$};
\end{axis}
\end{tikzpicture}
\end{codecote}
\begin{questions}[start=4]
  \item Résoudre sur $\intFF{-\pi}{\pi}$ l'équation $\cos x=\frac12$, puis l'inéquation $\cos x\leq\frac12$, à l'aide du cercle.
    Contrôler les résultats sur la courbe de la fonction cosinus tracée ci-dessus.
\end{questions}
\end{exercice}

\rubrique{Entraînement}

\begin{exercice}[2]{Parité et périodicité en images\enspace\otNouveau}
On considère les fonctions $f$, $g$, $h$ et $k$ définies sur $\R$ par
\[ f(x)=\cos(2x),\qquad g(x)=x\sin x,\qquad h(x)=\sin x+\cos x,\qquad k(x)=\sin x\cos x. \]
Leurs courbes sont tracées ci-dessous sur $\intFF{-2\pi}{2\pi}$, dans le désordre.\par\vspace{1mm}
\noindent\begin{tabular}{@{}c@{\hspace{3mm}}c@{}}
\otCourbeTrig{1}{1.3}{-1,-0.5,0.5,1}{sin(deg(x))*cos(deg(x))}{2.2} & \otCourbeTrig{2}{1.65}{-1,1}{sin(deg(x))+cos(deg(x))}{1.75} \\[1mm]
\otCourbeTrig{3}{5.4}{-4,-2,2,4}{x*sin(deg(x))}{0.55} & \otCourbeTrig{4}{1.3}{-1,-0.5,0.5,1}{cos(deg(2*x))}{2.2}
\end{tabular}
\begin{questions}
  \item Étudier la parité de chacune des quatre fonctions. Pour $h$, on calculera $h\left(\frac{\pi}{4}\right)$ et
    $h\left(-\frac{\pi}{4}\right)$.
  \item Montrer que $f$ et $k$ sont périodiques de période $\pi$, et que $h$ est périodique de période $2\pi$.
  \item Associer chaque fonction à sa courbe, en justifiant à l'aide des symétries et d'une valeur particulière.
  \item Expliquer pourquoi il suffit d'étudier $k$ sur $\intFF{0}{\frac{\pi}{2}}$ pour tracer toute sa courbe, et préciser les
    transformations à utiliser.
\end{questions}
\end{exercice}

\begin{exercice}[2]{Dérivées}
Calculer la dérivée de chacune des fonctions suivantes.
\begin{tasks}(3)
  \task $f(x)=\sin(3x)$
  \task $g(x)=x^{2}\cos x$
  \task $h(x)=\tan\left(x^{2}+1\right)$
  \task $k(x)=\sin x\cos x$
  \task $m(x)=\dfrac{\sin x}{x}$, sur $\intOO{0}{+\infty}$
\end{tasks}
Pour $k$, on exprimera le résultat à l'aide de $\cos(2x)$, sachant que $\cos(2x)=\cos^{2}x-\sin^{2}x$.
\end{exercice}

\begin{exercice}[2]{Équations}
Résoudre les équations suivantes sur l'intervalle $\intFO{0}{2\pi}$.
\begin{tasks}(3)
  \task $\sin x=\dfrac12$
  \task $\cos x=-\dfrac{\sqrt2}{2}$
  \task $\tan x=1$
  \task $2\sin^{2}x-3\sin x+1=0$
  \task $\sin(2x)=\cos x$
\end{tasks}
Pour la dernière, on utilisera $\sin(2x)=2\sin x\cos x$.
\end{exercice}

\begin{exercice}[2]{Inéquations}
\begin{questions}
  \item Résoudre $\sin x\geq\dfrac{\sqrt3}{2}$ sur $\intFF{0}{2\pi}$.
  \item Résoudre $\cos(2x)<\dfrac12$ sur $\intFF{0}{\pi}$.
  \item Résoudre $\tan x\leq1$ sur $\intFO{0}{2\pi}$, privé de $\frac{\pi}{2}$ et de $\frac{3\pi}{2}$.
  \item À l'aide de la formule $\sin(a+b)=\sin a\cos b+\cos a\sin b$, montrer que, pour tout réel $x$,
    $\sin x+\cos x=\sqrt2\sin\left(x+\frac{\pi}{4}\right)$. En déduire le maximum de la fonction $x\mapsto\sin x+\cos x$
    sur $\intFF{0}{2\pi}$, puis résoudre l'inéquation $\sin x+\cos x\geq1$ sur $\intFF{0}{2\pi}$.
\end{questions}
\end{exercice}

\begin{exercice}[2]{Formules et identités}
On rappelle que, pour tous réels $a$, $b$ et $x$ :
\[ \cos(a-b)=\cos a\cos b+\sin a\sin b \qquad \sin(a\pm b)=\sin a\cos b\pm\cos a\sin b \]
\[ \cos(2x)=2\cos^{2}x-1=1-2\sin^{2}x \qquad \sin(2x)=2\sin x\cos x \]
\begin{questions}
  \item Simplifier $\sin(a+b)+\sin(a-b)$.
  \item Calculer $\cos\dfrac{\pi}{12}$, en remarquant que $\dfrac{\pi}{12}=\dfrac{\pi}{3}-\dfrac{\pi}{4}$.
  \item Démontrer que, pour tout réel $x$, $\sin^{2}x=\dfrac{1-\cos(2x)}{2}$.
  \item Simplifier $\dfrac{\sin(2x)}{1+\cos(2x)}$ pour tout réel $x$ tel que $\cos x\neq0$.
  \item Calculer $\sin\dfrac{5\pi}{12}$.
\end{questions}
\end{exercice}

\begin{exercice}[2]{Limites}
\begin{questions}
  \item Rappeler la valeur de $\displaystyle\lim_{x\to0}\frac{\sin x}{x}$.
  \item En multipliant numérateur et dénominateur par $1+\cos x$, calculer $\displaystyle\lim_{x\to0}\frac{1-\cos x}{x^{2}}$.
  \item Calculer $\displaystyle\lim_{x\to\frac{\pi}{2}}\frac{\tan x}{x-\frac{\pi}{2}}$, en distinguant la limite à gauche et la limite à droite.
  \item La fonction $f:x\mapsto\dfrac{\sin x}{x}$, définie sur $\R^{*}$, peut-elle être prolongée par continuité en $0$ ?
  \item Calculer $\displaystyle\lim_{x\to+\infty}\frac{\sin x}{x}$.
\end{questions}
\end{exercice}

\rubrique{Problèmes}

\begin{exercice}[3]{Études de fonctions}
Étudier les variations de chacune des fonctions suivantes sur l'intervalle indiqué : dérivée, signe de la dérivée,
tableau de variations et extremums.
\begin{questions}
  \item $f(x)=\sin x+\cos x$ sur $\intFF{0}{2\pi}$.
  \item $g(x)=\sin^{2}x-\cos^{2}x$ sur $\intFF{0}{2\pi}$. On pourra utiliser $\cos(2x)=\cos^{2}x-\sin^{2}x$.
  \item $h(x)=\dfrac{\sin x}{1+\cos x}$ sur $\intFO{0}{2\pi}$ privé de $\pi$. Préciser les limites en $\pi$.
  \item $k(x)=\sin x\cos x$ sur $\intFF{0}{2\pi}$. On pourra utiliser $\sin(2x)=2\sin x\cos x$.
  \item $m(x)=\tan x-x$ sur $\intFO{0}{\frac{\pi}{2}}$. En déduire que $\tan x\geq x$ pour tout $x\in\intFO{0}{\frac{\pi}{2}}$.
\end{questions}
\end{exercice}

\begin{exercice}[3]{Synthèse}
\begin{questions}
  \item Résoudre l'équation $\sin(3x)=\sin x$ sur $\intFF{0}{2\pi}$.
  \item Étudier les variations de $f(x)=\dfrac{\sin x}{1+\sin x}$ sur $\intFF{0}{2\pi}$ privé de $\frac{3\pi}{2}$, et ses limites en $\frac{3\pi}{2}$.
  \item Résoudre l'inéquation $\cos(2x)\geq\dfrac12$ sur $\intFF{0}{2\pi}$.
  \item Calculer la dérivée seconde de la fonction $g$ définie sur $\R$ par $g(x)=\sin\left(x^{2}\right)$.
  \item Résoudre l'équation $\tan(2x)=\sqrt3$ sur $\intFF{0}{\pi}$.
\end{questions}
\end{exercice}

\begin{exercice}[3]{Suites et trigonométrie}
\begin{questions}
  \item On pose $u_n=\cos\left(\dfrac{\pi}{2^{n}}\right)$ pour $n\in\mathbb{N}$. Calculer $u_0$, $u_1$, $u_2$, puis $\displaystyle\lim_{n\to+\infty}u_n$.
  \item On pose $v_n=\sin\left(\dfrac{n\pi}{4}\right)$. Montrer que la suite $(v_n)$ est périodique et déterminer sa période.
  \item On pose $w_n=\dfrac1n\displaystyle\sum_{k=1}^{n}\sin\left(\frac{k\pi}{n}\right)$ pour $n\geq1$. Calculer $w_2$ et $w_4$.
  \item La suite $(u_n)$ est définie par $u_0=2$, $u_1=2\cos\frac{\pi}{3}$ et, pour tout $n\geq1$,
    $u_{n+1}=2\cos\left(\frac{\pi}{3}\right)u_n-u_{n-1}$. Calculer $u_1$ et $u_2$, puis démontrer par récurrence que
    $u_n=2\cos\left(\dfrac{n\pi}{3}\right)$ pour tout $n\in\mathbb{N}$. On pourra utiliser
    $2\cos a\cos b=\cos(a+b)+\cos(a-b)$.
  \item On pose $t_n=\displaystyle\sum_{k=0}^{n}\cos\left(\frac{k\pi}{3}\right)$. Calculer $t_1$, $t_2$ et $t_6$, puis montrer
    que la suite $(t_n)$ est périodique.
\end{questions}
\end{exercice}

\begin{exercice}[3]{Étude complète d'une fonction trigonométrique}
On considère la fonction $f$ définie sur $\R$ par $f(x)=2\sin x+\sin(2x)$, de courbe $\Cf$.
\partieExo{A}{Réduction de l'intervalle d'étude}
\begin{questions}
  \item Montrer que $f$ est périodique de période $2\pi$ et qu'elle est impaire. Quelles transformations du plan laissent
    $\Cf$ invariante ?
  \item Expliquer pourquoi il suffit d'étudier $f$ sur $\intFF{0}{\pi}$.
\end{questions}
\partieExo{B}{Variations}
\begin{questions}[resume]
  \item On admet que $\cos(2x)=2\cos^{2}x-1$ pour tout réel $x$. Montrer que, pour tout réel $x$,
    $f'(x)=2(2\cos x-1)(\cos x+1)$.
  \item Étudier le signe de $f'(x)$ sur $\intFF{0}{\pi}$, puis dresser le tableau de variations de $f$ sur $\intFF{0}{\pi}$.
    Donner la valeur exacte du maximum.
\end{questions}
\partieExo{C}{Une équation et une tangente}
\begin{questions}[resume]
  \item Vérifier que $\frac{\pi}{2}$ est solution de l'équation $f(x)=2$, puis montrer que cette équation admet exactement deux
    solutions sur $\intFF{0}{\pi}$. Donner un encadrement d'amplitude $10^{-2}$ de celle qui est différente de $\frac{\pi}{2}$, notée $\alpha$.
  \item Déterminer une équation de la tangente $T$ à $\Cf$ à l'origine. Montrer que, pour tout réel $x$,
    $f'(x)-4=2\left(\cos x-1\right)\left(2\cos x+3\right)$, puis en déduire la position de $\Cf$ par rapport à $T$ sur $\intFO{0}{+\infty}$.
  \item Tracer $\Cf$ sur $\intFF{-\pi}{\pi}$, ainsi que la droite $T$.
\end{questions}
\end{exercice}

\begin{exercice}[3]{Oscillations amorties}
On écarte un ressort de sa position d'équilibre, puis on le lâche. La position de son extrémité, en centimètres, à
l'instant $t\geq0$ exprimé en secondes, est donnée par $x(t)=e^{-t/2}\cos(2t)$.
\partieExo{A}{Enveloppes et limite}
\begin{questions}
  \item Montrer que, pour tout $t\geq0$, $-e^{-t/2}\leq x(t)\leq e^{-t/2}$. En déduire la limite de $x(t)$ lorsque $t$ tend vers $+\infty$, et interpréter.
  \item Déterminer les instants auxquels le ressort passe par sa position d'équilibre, c'est-à-dire les solutions de $x(t)=0$.
  \item Résoudre l'équation $x(t)=e^{-t/2}$ sur $\intFO{0}{+\infty}$.
\end{questions}
\partieExo{B}{Vitesse et amortissement}
\begin{questions}[resume]
  \item Montrer que, pour tout $t\geq0$, $x'(t)=-e^{-t/2}\left(\frac12\cos(2t)+2\sin(2t)\right)$. Montrer qu'aux instants trouvés
    à la question 3, la courbe de $x$ et celle de $t\mapsto e^{-t/2}$ ont la même tangente.
  \item Montrer que, pour tout $t\geq0$, $x(t+\pi)=e^{-\pi/2}\,x(t)$. Interpréter ce résultat.
  \item On considère que les oscillations deviennent imperceptibles dès que $e^{-t/2}<0{,}01$. À partir de quel instant est-ce le cas ?
\end{questions}
\end{exercice}

\begin{exercice}[3]{Encadrements de sinus et cosinus}
\partieExo{A}{Trois inégalités}
Dans cette partie, $x$ désigne un réel positif.
\begin{questions}
  \item Étudier les variations de la fonction $g:x\mapsto x-\sin x$ sur $\intFO{0}{+\infty}$, puis en déduire que $\sin x\leq x$.
  \item On pose $h(x)=\cos x-1+\dfrac{x^{2}}{2}$. Calculer $h'(x)$ et en déduire, à l'aide de la question 1, que
    $1-\dfrac{x^{2}}{2}\leq\cos x$.
  \item De la même manière, démontrer que $x-\dfrac{x^{3}}{6}\leq\sin x$.
\end{questions}
\partieExo{B}{Applications}
\begin{questions}[resume]
  \item En déduire un encadrement de $\sin(0{,}1)$ d'amplitude inférieure à $2\times10^{-4}$, ainsi qu'une minoration de $\cos(0{,}2)$.
  \item Montrer que, pour tout $x>0$, $1-\dfrac{x^{2}}{6}\leq\dfrac{\sin x}{x}\leq1$. Retrouver ainsi la limite de $\dfrac{\sin x}{x}$
    lorsque $x$ tend vers $0$.
\end{questions}
\end{exercice}
